\documentclass[10pt]{article}
\usepackage[a4paper,margin=0.75in]{geometry}
\usepackage{times}
\usepackage{microtype}
\usepackage{booktabs}
\usepackage{array}
\usepackage{amsmath}
\usepackage{enumitem}
\usepackage{xcolor}
\usepackage{titlesec}
\usepackage[hidelinks]{hyperref}

\setlist{nosep,leftmargin=1.2em}
\titlespacing*{\section}{0pt}{6pt}{2pt}
\titlespacing*{\subsection}{0pt}{4pt}{1pt}
\titleformat{\section}{\normalfont\large\bfseries}{\thesection}{0.5em}{}
\titleformat{\subsection}{\normalfont\normalsize\bfseries}{\thesubsection}{0.5em}{}
\setlength{\parindent}{0pt}
\setlength{\parskip}{3pt}

% Placeholder for numbers the team fills in after the GPU runs.
\newcommand{\fic}{\textcolor{gray}{$\cdot$}} % fill-in cell

\title{\vspace{-2.5em}\textbf{Generative Retrieval for Recommendation:\\
a Reproduction and Controlled Study of TIGER}}
\author{Anna Kiseleva \quad Maya Lager \quad Kseniia Tsypliakova\\
\normalsize HSE University --- DeepRecSys final project}
\date{\vspace{-1em}}

\begin{document}
\maketitle
\vspace{-1.5em}

\begin{abstract}
\noindent
TIGER (Rajput et al., NeurIPS 2023) reframes sequential recommendation as
\emph{generative retrieval}: each item is mapped to a short tuple of discrete
``Semantic ID'' tokens via an RQ-VAE over content embeddings, and a seq2seq
Transformer is trained from scratch to autoregressively generate the next item's
Semantic ID. We reimplement the full pipeline from scratch (no original code is
public) and build a controlled study of the design choices the paper leaves
implicit. The code reproduces the TIGER and SASRec pipelines on Amazon Beauty
under a strict full-ranking leave-one-out protocol, then asks four questions:
whether residual quantization is necessary (vs a prefix-conditional hierarchical
$k$-means tree), how Semantic ID
\emph{structure} (token order, collision-token placement, code length) affects
quality and failure modes, whether explicit time-gap tokens help, and whether the
autoregressive factorization is necessary at all (vs a non-autoregressive
masked-denoising decoder). Beyond accuracy we report generative-retrieval
diagnostics --- collision rate, invalid-code rate, and decode latency --- that
accuracy-only tables hide.
\end{abstract}

\vspace{-0.5em}
\section{Experimental setup}

\textbf{Data.} Amazon Product Reviews (2014), category \emph{Beauty} as the
primary benchmark, with \emph{Sports and Outdoors} and \emph{Toys and Games} as
optional replications. For the paper-style leave-one-out reproduction we use the
standard full-data 5-core file; Beauty then has $22{,}363$ users and $12{,}101$
items (mean sequence length $8.9$), matching the paper. For temporal
experiments, we instead preprocess the full review file and fit the iterative
k-core filter only on the train window ($t< t_{\mathrm{val}}$), then keep later
interactions for validation/test. This prevents future interactions from
deciding which users/items enter the temporal dataset. Items are described by a
sentence built from title, brand, price, and categories, embedded with a frozen
Sentence-T5 encoder ($768$-dim), matching the paper.

\textbf{Evaluation protocol.} We use two splits: leave-one-out (last interaction
= test, second-to-last = validation, rest = training with all-prefix
augmentation), used to reproduce the paper; and a global \emph{temporal} split
(two timestamp cutoffs; history always precedes the target in wall-clock time),
used for the quality and time-aware experiments, since leave-one-out is a weak
quality protocol. Crucially we evaluate against the \textbf{full item catalogue,
not sampled negatives} --- sampled metrics are a known source of inflated,
non-comparable numbers. We report Recall@$K$ and NDCG@$K$ for $K\in\{5,10\}$.
For temporal experiments the user/item core, Semantic-ID quantizer/tree, and
collision-token priority are all determined from the temporal training window and
then used to assign codes to the whole catalogue; otherwise future interactions
or future item content would leak into the protocol.
Collision/disambiguation tokens are also assigned to train-window items first,
so future or cold catalogue items with the same content code cannot shift a
training item's target Semantic ID.
For generative models the candidate set is what beam search (beam $=100$)
decodes; the ground-truth rank is read from that ordered list and items never
generated count as misses. Headline TIGER and NAR numbers use
\textbf{unconstrained} decoding (as in the paper: invalid generations are
dropped and counted as a failure mode); trie-constrained decoding is reported
only as an ablation, since it cannot emit invalid codes and is therefore
optimistic. All trainers select the checkpoint with the best Recall@10 on the
full validation split; the optional evaluation-user cap is used only for
development runs. For temporal splits we also report the fraction of
validation/test targets already seen in the training window, because cold target
items affect SASRec and content-code models differently.

\textbf{Models.} \emph{SASRec} (self-attentive next-item, full-softmax) is the
main baseline. \emph{TIGER} is a T5 encoder--decoder ($d{=}128$, 4+4 layers,
trained from scratch) over the Semantic ID token space, with per-position
codebook offsets, a user-hash token, and a trie for prefix-constrained decoding.
Semantic IDs use $3$ RQ-VAE levels with codebook size $256$ plus one collision
token, giving $4$-token codes --- as in the paper.

\textbf{Reproducibility.} A single \texttt{scripts/run\_beauty.sh} runs the whole
pipeline (download $\rightarrow$ preprocess $\rightarrow$ Semantic IDs
$\rightarrow$ all experiments). Code, README, and a synthetic smoke test that
exercises the pipeline end-to-end without any downloads are included.

\section{Research questions}
\begin{description}[leftmargin=2.2em,style=nextline]
\item[\textmd{\textbf{RQ1}}] \emph{Is residual quantization necessary?} Do RQ-VAE
Semantic IDs beat a simpler prefix-conditional content tree (hierarchical
$k$-means) built from the \emph{same} embeddings and code shape?
\item[\textmd{\textbf{RQ2}}] \emph{How does Semantic ID structure affect quality
and failure modes?} We vary token order (original / reversed / permuted),
collision-token placement (last vs first), and code length / codebook size.
\item[\textmd{\textbf{RQ3}}] \emph{Do explicit time signals help?} We discretize
inter-event gaps into 5 buckets and compare two ways of injecting them: a
separate gap \emph{token} interleaved between items, and --- following the course
feedback --- a gap \emph{embedding} added to each event's token embeddings
(analogous to SASRec's positional embedding), evaluated under a temporal split.
\item[\textmd{\textbf{RQ4}}] \emph{Is the autoregressive pipeline necessary?} We
compare against two non-autoregressive decoders: a MaskGIT-style masked-denoising
model, and a simpler baseline that predicts all code positions in parallel with
independent linear heads. Their headline metrics use unconstrained tuple
enumeration, with trie-constrained enumeration kept as an ablation.
\end{description}

\section{Results and discussion}

\textbf{Reproduction (RQ0).} Table~\ref{tab:main} shows the paper's reported
Beauty numbers next to our runs. The current table is intentionally left as a
fill-in table until the full GPU runs are complete. Because TIGER is trained from
scratch on a single category, run-to-run variance is high; we will treat ``same
ballpark and same ordering (TIGER $>$ SASRec on NDCG)'' as a successful
reproduction rather than digit-exact matching.

\begin{table}[h]
\centering
\small
\caption{Amazon Beauty, full-ranking leave-one-out. Paper values from Rajput et
al. (2023, Table 1). ``Ours'' = our reimplementation.}
\label{tab:main}
\begin{tabular}{l cccc}
\toprule
& Recall@5 & NDCG@5 & Recall@10 & NDCG@10 \\
\midrule
SASRec (paper) & 0.0387 & 0.0249 & 0.0605 & 0.0318 \\
TIGER (paper)  & 0.0454 & 0.0321 & 0.0648 & 0.0384 \\
\midrule
SASRec (ours)  & \fic & \fic & \fic & \fic \\
TIGER (ours)   & \fic & \fic & \fic & \fic \\
\bottomrule
\end{tabular}
\end{table}

\textbf{Ablations (RQ1--RQ4).} Table~\ref{tab:abl} collects the controlled
variants, all trained and evaluated identically and compared against our own
TIGER reproduction (the ``TIGER (ours)'' row of Table~\ref{tab:main}). Each row
changes exactly one design choice. We also log diagnostics that pure accuracy
hides: the pre-disambiguation \emph{collision rate} (how many items share a
content code, i.e.\ how much work the non-semantic collision token is doing), the
\emph{invalid-code rate} under \emph{unconstrained} decoding (overall and at
$K\in\{5,10,20\}$; constrained decoding forces this to $\approx0$), mean
\emph{decode latency} per user, valid-candidate coverage after invalid/duplicate
filtering, and target-coverage diagnostics for temporal splits.

\begin{table}[h]
\centering
\small
\caption{Controlled variants on Beauty ($\Delta$ vs our TIGER reproduction).
Fill from \texttt{runs/*.json}.}
\label{tab:abl}
\begin{tabular}{l l cc c}
\toprule
RQ & Variant & R@10 & N@10 & Key diagnostic \\
\midrule
-- & TIGER (RQ-VAE, original) & \fic & \fic & coll.\ rate \fic \\
RQ1 & Hierarchical $k$-means IDs & \fic & \fic & coll.\ rate \fic \\
RQ2 & Token order: reversed & \fic & \fic & \\
RQ2 & Token order: permuted & \fic & \fic & \\
RQ2 & Collision token first & \fic & \fic & \\
RQ2 & Code length $L{=}4$ levels & \fic & \fic & coll.\ rate \fic \\
RQ3 & + time-gap token & \fic & \fic & \\
RQ3 & + time-gap embedding & \fic & \fic & temporal split \\
RQ4 & Masked-denoising (NAR) & \fic & \fic & inv.\ code \fic; lat.\ \fic \\
RQ4 & Parallel linear heads & \fic & \fic & inv.\ code \fic; lat.\ \fic \\
\bottomrule
\end{tabular}
\end{table}

\textbf{What we expect to find, and how we will read it.}
\begin{itemize}
\item \textbf{RQ1.} If RQ-VAE $\approx$ hierarchical $k$-means on accuracy, the
``residual'' machinery may not be the only driver of TIGER --- a simpler
tree-structured content discretizer may suffice. This baseline is not a perfect
drop-in replacement for residual quantization because its token meanings are
prefix-conditional, so the conclusion is framed as an empirical comparison
against a simpler content-code tree. A gap in favour of RQ-VAE supports the
paper's hierarchy claim. We also compare collision rates: a method that collides
less leans less on the non-semantic disambiguation token.
\item \textbf{RQ2.} Reversing/permuting token order tests whether the model
exploits the coarse-to-fine hierarchy (a large drop $\Rightarrow$ order carries
real semantic signal). Moving the collision token to the front tests whether a
non-semantic token early in the code disrupts decoding. The code-length sweep
trades off collision rate (longer codes $\Rightarrow$ fewer collisions) against a
harder generation target.
\item \textbf{RQ3.} We contrast two ways of adding the same gap information: a
separate token vs an embedding summed into each event. If the embedding helps
where the token does not, the signal is useful but the token form wastes
sequence length and dilutes attention; if neither helps, the autoregressive
context already captures order. We read these under the temporal split, where
recency genuinely matters.
\item \textbf{RQ4.} Both non-autoregressive decoders are much faster (one forward
pass vs $L$ steps) but cannot model left-to-right dependencies between code
positions; the parallel-head model additionally drops all cross-position
modelling. The accuracy ladder autoregressive $\to$ masked-denoising $\to$
parallel heads, read together with each model's \emph{unconstrained} invalid-code
rate and latency, quantifies how much the autoregressive factorization actually
buys --- the core question of whether generative retrieval \emph{needs} to be
autoregressive.
\end{itemize}

\textbf{Methodology and honest caveats.} This is a scaled-down reimplementation,
not a bit-exact reproduction: our RQ-VAE uses a far smaller training budget than
the paper (500 epochs, Adam, vs $\sim$20k steps, Adagrad), so we target the right
ballpark and the right \emph{ordering}, not identical digits. We verified the
generative evaluation end-to-end with unit tests (notably a Semantic-ID
$\leftrightarrow$ item round-trip test after we found and fixed an off-by-one in
the code$\to$item map that had been corrupting the generative metrics), and we
report TIGER under paper-style unconstrained decoding with constrained decoding
kept as a clearly-labelled ablation. The same headline-vs-ablation decoding
distinction is used for NAR, so RQ4 does not conflate autoregressive
factorization with trie constraints.

\section{Conclusion}
We provide a clean, fully reproducible reimplementation of TIGER and a controlled
study that isolates four design choices the original paper bundles together,
evaluated under full-ranking protocols and augmented with
generative-retrieval diagnostics. The framework lets each research question be
answered by a single command and a single changed flag, making both the positive
and the negative results interpretable. Final numbers and per-RQ conclusions are
populated from the GPU runs (\texttt{runs/*.json}); the analysis template above
states in advance how each outcome should be read, so the conclusions follow
directly from the measured deltas.

\vspace{0.3em}
{\footnotesize \textbf{References.} S.\ Rajput et al., ``Recommender Systems with
Generative Retrieval,'' NeurIPS 2023 (arXiv:2305.05065).\quad W.-C.\ Kang, J.\
McAuley, ``Self-Attentive Sequential Recommendation,'' ICDM 2018.}

\end{document}
